\section{Constructive integer projection without residue tables}
\label{sec:projection}
\status{Implemented in \code{presburger.py}; exact symbolic witness programs,
feasibility guards, Bezout receipts, and differential checks.}

\subsection{The exact supported fragment}
Let $x=(x_1,\ldots,x_n)\in\Z^n$ be parameters and let $y\in\Z$ be the output.
The input is a conjunction
\begin{equation}
\label{eq:projection-input}
 \Phi(x,y)\equiv
 \bigwedge_i a_i y\le\beta_i(x)
 \quad\land\quad
 \bigwedge_j c_jy\equiv\delta_j(x)\pmod{m_j},
\end{equation}
where the coefficients $a_i,c_j,m_j$ are fixed integers, $m_j>0$, and the right
sides are integer affine expressions in $x$. Coefficients may be negative or
zero. An empty conjunction is allowed. There is exactly one existential output
variable in this initial fragment.

The worker returns an integer expression $w(x)$ and a quantifier-free guard
$C(x)$ such that
\begin{equation}
\label{eq:projection-contract}
 C(x)\ \Longleftrightarrow\ \exists y\in\Z,\Phi(x,y),
 \qquad C(x)\ \Longrightarrow\ \Phi(x,w(x)).
\end{equation}
The expression uses affine operations, positive-denominator floor division,
remainder, minimum, and maximum. It is a single program for all parameters, not
a table of unrelated witnesses for sampled inputs.

This directly extends the useful area between Forge's integer-lattice worker
and its finite-residue witness tables. A finite table can cost proportional to a
modulus whose binary representation is short. The construction here uses the
arithmetic structure of the congruences instead.

\subsection{Normalize arbitrary fixed coefficients in congruences}
For one congruence
\[
 c y\equiv\delta\pmod m,
\]
let $g=\gcd(c,m)>0$ and choose integers $u,v$ with $uc+vm=g$. A solution exists
only if $g\mid\delta$. Under that condition the congruence is equivalent to
\begin{equation}
\label{eq:normalize-congruence}
 y\equiv u\frac{\delta}{g}\pmod{m/g}.
\end{equation}
The quotient is an exact integer quotient under the guard. In the program it
is computed with floor division, but the divisibility guard is never discarded.
When $c=0$, the normalization reduces to the guard $m\mid\delta$ and a modulus
of one, imposing no restriction on $y$.

The certificate does not ask the checker to run the extended Euclidean
algorithm. It supplies $(g,u,v)$, and the checker verifies
\[
 g>0,\quad g\mid c,\quad g\mid m,\quad uc+vm=g.
\]
These conditions imply that $g$ is the positive gcd: every common divisor divides
the displayed integer combination, while $g$ itself is a common divisor.

\subsection{Merge congruences with non-coprime moduli}
Suppose the current system has been reduced to $y\equiv r\pmod M$, and the next
normalized congruence is $y\equiv s\pmod N$, with $M,N>0$. Put
\[
 h=\gcd(M,N),\qquad pM+qN=h.
\]
Compatibility requires $h\mid(s-r)$. Under that guard, set
\begin{align}
 k&=\left(p\frac{s-r}{h}\right)\bmod(N/h),\label{eq:crt-k}\\
 M'&=M(N/h)=\lc(M,N),\nonumber\\
 r'&=(r+Mk)\bmod M'.\label{eq:crt-r}
\end{align}
The two congruences are equivalent to $y\equiv r'\pmod{M'}$. In the total program, the displayed quotient $(s-r)/h$ is floor division; compatibility makes it exact. The residue expression therefore remains defined even when the compatibility guard is false.

\begin{lemma}[Guarded CRT merge]
\label{lem:crt}
For the definitions above,
\[
 (y\equiv r\pmod M\ \land\ y\equiv s\pmod N)
 \Longleftrightarrow
 (h\mid(s-r)\ \land\ y\equiv r'\pmod{M'}).
\]
\end{lemma}
\begin{proof}
If both congruences hold, writing $y-r=M\ell$ and $y-s=Nj$ shows
$s-r=M\ell-Nj$, hence compatibility. Dividing by $h$, the equation for
$\ell$ modulo $N/h$ has inverse $p$ for $M/h$, because
$p(M/h)+q(N/h)=1$. Therefore $\ell\equiv k\pmod{N/h}$, which is equivalent
to $y\equiv r+Mk\pmod{M'}$. Conversely, compatibility and that last
congruence give the first congruence directly and the second by substituting
\eqref{eq:crt-k}. Reducing $r+Mk$ modulo $M'$ does not change its class.
\end{proof}

Start with residue $r=0$ and period $M=1$, then merge every congruence. This
produces a fixed final period and a symbolic residue program $r(x)$, together
with divisibility guards. The moduli need not be prime, pairwise coprime, or
small. Every normalized modulus and every merge uses its own checked Bezout
receipt.

\subsection{Convert inequalities into an interval}
For $a y\le\beta$:
\[
 \begin{array}{rcll}
 a>0 &\Rightarrow& y\le\lfloor\beta/a\rfloor, &\text{an upper bound},\\[2pt]
 a=-c<0 &\Rightarrow& y\ge-\lfloor\beta/c\rfloor, &\text{a lower bound},\\[2pt]
 a=0 &\Rightarrow& 0\le\beta, &\text{a parameter guard}.
 \end{array}
\]
The negative-coefficient case is easy to get wrong. It uses
$\lceil-\beta/c\rceil=-\lfloor\beta/c\rfloor$, with $c>0$. Truncation toward
zero would give incorrect answers for negative parameters.

Let $L(x)$ be the maximum of the lower bounds when any exist, and $U(x)$ the
minimum of the upper bounds when any exist. Do not invent sentinel values for
missing endpoints: the problem is over unbounded integers, not a machine range.

\subsection{Choose the nearest congruent point}
After normalization, all solutions to the congruences have the form $r+Mz$.
If a lower bound exists, define
\begin{equation}
\label{eq:lower-witness}
 w=L+((r-L)\bmod M).
\end{equation}
This is the smallest integer at least $L$ in the required residue class. If an
upper bound also exists, add the final guard $w\le U$.

If there is no lower bound but there is an upper bound, define instead
\begin{equation}
\label{eq:upper-witness}
 w=U-((U-r)\bmod M).
\end{equation}
This is the largest congruent integer at most $U$, so no extra interval
feasibility guard is required. With neither bound, use $w=r$.

\begin{theorem}[Constructive projection]
\label{thm:projection}
The conjunction of the normalization guards, CRT compatibility guards,
zero-coefficient inequality guards, and the optional $w\le U$ guard satisfies
\eqref{eq:projection-contract}, with the witness defined above.
\end{theorem}
\begin{proof}
The normalization and CRT lemmas establish an equivalence between the original
congruences and their guards together with $y\equiv r\pmod M$. The inequality
transformations are equivalences over integers, including the zero-coefficient
cases. It remains to intersect one arithmetic progression with one possibly
unbounded interval.

With a lower endpoint, \eqref{eq:lower-witness} is congruent to $r$ and lies in
$[L,L+M)$. Any other congruent integer at least $L$ differs from it by a
nonnegative multiple of $M$. Thus an upper endpoint can be satisfied precisely
when $w\le U$. With only an upper endpoint, \eqref{eq:upper-witness} always
provides a suitable congruent point. With no endpoints, $r$ itself suffices.
These arguments prove both feasibility equivalence and correctness of the
constructed witness, for every parameter valuation.
\end{proof}

\subsection{A witness program, not an exponentially large normal form}
The implementation stores expressions as a hash-consed directed acyclic graph.
A node refers only to earlier nodes and is one of:
\begin{lstlisting}
const(z)        var(i)          scale(z, node)
add(a, b)       div(a, d)       mod(a, d)     // d > 0
max(a, b)       min(a, b)
\end{lstlisting}
Guards are divisibility tests or comparisons of node values. There is no floating
point operation, no variable denominator, and no implicit conversion to machine
integers. Shared affine expressions and intermediate residues need not be
copied into a large tree-shaped term.

For fixed parameter arity, program-node count grows linearly with the number of
constraints, apart from sharing effects. More explicitly, assembling affine
right sides costs $O((p+q)(n+1))$ nodes for $p$ inequalities and $q$
congruences. Integer coefficient bit lengths and the cost of division still
matter. The final period can be a product of many moduli, but its bit length is
at most the sum of their bit lengths before normalization. No loop runs once
per residue class.

\begin{example}[A 60-bit period in 42 nodes]
The recorded large case has
\[
 m_1=1{,}000{,}000{,}007,\qquad m_2=1{,}000{,}000{,}009,
\]
\[
 M=m_1m_2=1{,}000{,}000{,}016{,}000{,}000{,}063,
\]
and constraints
\[
 x\le y\le x+M-1,\qquad
 y\equiv2x+1\pmod{m_1},\qquad
 y\equiv5x-3\pmod{m_2}.
\]
The generated program has 42 nodes. The final interval has exactly one
representative of every residue class, so the theorem establishes feasibility
for all $x$. At $x=0$ the generated witness is $2{,}000{,}000{,}015$. The
experiment also checks the program concretely at $x=\pm10^{100}$. Those large
values are implementation tests; the universal claim follows from the
construction theorem, not from their size or number.
\end{example}

A residue table for this combined period is not a sensible representation.
This is a representation-size contrast, not a measured speed comparison with
Forge's existing residue implementation.

\subsection{How to discharge a quantified Lean goal}
For a goal
\[
 \forall x,\quad D(x)\longrightarrow\exists y:\Z,\Phi(x,y),
\]
introduce $x$ and the domain assumptions, generate $(C,w)$, and create the
remaining obligation $D(x)\Rightarrow C(x)$. Once this is proved, the generic
projection theorem supplies the existential with witness $w(x)$. A local
arithmetic solver can attempt the guard, while another worker may supply a
missing congruence or bound. There is no reason to enumerate witnesses after
the exact projection has already constructed one.

Conversely, a valuation of $x$ violating $C$ is useful only after checking the
domain $D$ and replaying the projection equivalence. A raw solver model for an
unrelated abstraction is not evidence that the original universally quantified
goal is false.

For natural-number output, first add $0\le y$ and prove that the selected
integer can be transported to \code{Nat}. For rational or real parameters,
this integer theorem is not the correct source semantics. Existential
quantifier order also matters: a witness may depend only on parameters in
scope before its binder.

\subsection{The remaining verification work}
The Python checker validates the Bezout receipts and deterministically
reconstructs the witness program and guards from the input. It then compares
this reconstruction with the supplied program. Search and checking share that
assembler. Thus replay detects modified receipts or programs but cannot,
by itself, detect an arithmetic mistake common to both uses of the assembler.
The direct constraint evaluator and exhaustive concrete oracle in the test suite
address that risk by a different route; they do not amount to a formal proof.

The Lean implementation should prove normalization, guarded CRT, bound
conversion, and nearest-congruent-point lemmas once. Certificate elaboration
then composes these lemmas, with positive-divisor facts explicit. A reflected
program interpreter is an optional later optimization; an initial version can
emit a shared \code{let}-term and ordinary proof applications. \code{omega}
is an appropriate solver to try on residual linear integer obligations, but
acceptance of this exact generated syntax must be measured rather than assumed.
The mathematical construction does not require any SMT solver.

\subsection{Limits that must survive the interface}
Variable coefficients such as $x\,y\le z$, variable moduli, nonlinear parameter
right sides in the current input grammar, general Boolean combinations, and
several mutually constrained existential outputs are outside the implemented
fragment. Disequalities can be handled by explicit branching in a future layer,
but they are not silently dropped. Repeatedly projecting several variables may
cause large piecewise expressions and requires a separate supported-language
argument after each elimination. A one-output completeness theorem must not be
advertised as complete arbitrary Presburger synthesis.
